\chapter{Maximum Mass for Macroscopic Quantum Superposition in Fractal T0 Geometry}


	
	\subsection*{Brief Introduction}
	
	This chapter examines how the fractal vacuum structure of FFGFT modifies gravitational decoherence of macroscopic quantum superpositions. In standard physics, several models for gravitational decoherence exist (Diósi-Penrose, Károlyházy), which are heuristically motivated. FFGFT offers a structural mechanism through the fractal nonlinearity of the vacuum field.
	
	\subsection*{Mathematical Foundation}
	
	The decoherence mechanism arises from the fractal nonlinearity of the vacuum field \(\Phi = \rho(x,t) e^{i\theta(x,t)}\). The modification is a structural consequence of the parameter \(\xi = \frac{4}{3} \times 10^{-4}\).
	

	
	\subsection*{Decoherence through the Fractal Vacuum}
	
	An object in superposition with spatial separation \(\Delta x\) causes different local gravitational fields in the two branches. This difference \(\Delta g\) leads to different phase gradients in the vacuum:
	
	\begin{equation}
		\Delta g = \frac{G M}{(\Delta x/2)^2} \cdot 2 \approx \frac{8 G M}{(\Delta x)^2}.
	\end{equation}
	
	The formula computes the acceleration difference between the two positions of the object — each superposition branch generates a gravitational field that acts at the location of the other branch.
	
	In FFGFT, the fractal phase gradient couples to \(\Delta g\). The decoherence rate acquires an additional \(\xi\)-dependent contribution:
	
	\begin{equation}
		\Gamma_{\text{FFGFT}} = \Gamma_{\text{grav}} \cdot g(\xi, \Delta x),
	\end{equation}
	
	where \(\Gamma_{\text{grav}}\) is the gravitational decoherence rate (e.g.\ Diósi-Penrose) and \(g(\xi, \Delta x)\) is a function of the fractal geometry describing the coupling between phase gradient and gravitational field difference.
	
	\textbf{Unit check:}
     \[
\begin{aligned}
	[\Gamma]
	&= \text{dimensionless}
	\cdot \si{\meter\cubed\per\kilo\gram\per\second\squared} \\
	&\quad \cdot
	\frac{\si{\kilo\gram}}{
		\si{\meter\per\second\squared} \cdot \si{\meter^2}
	} \\
	&= \si{\per\second}.
\end{aligned}
\]
	
	\subsection*{Qualitative Prediction}
	
	\textbf{Important caveat:} The technical reference documents (Doc~201) list ``modified decoherence rates in isolated systems'' and ``potential gravitational decoherence from T0 mediator'' as testable predictions, but do not give an explicit formula for a maximum mass \(M_{\max}\). The following analysis presents the \emph{structural mechanism} without claiming a specific numerical value.
	
	The FFGFT contribution to decoherence has the following properties:
	
	\begin{itemize}
		\item \textbf{Additional source:} The fractal nonlinearity of the vacuum generates decoherence \emph{beyond} standard environmental decoherence
		\item \textbf{ξ-scaling:} The contribution scales with \(\xi\), making it small compared to environmental effects but detectable in principle
		\item \textbf{Separation dependence:} The fractal correlation function \(f(\Delta x)\) introduces logarithmic corrections that differ from Diósi-Penrose
	\end{itemize}
	
	For the absolute scale of a maximum mass, a complete calculation of the fractal correlation function would be needed, which depends on parameters not fully specified in the current reference documents.
	
	\subsection*{Comparison with the Diósi-Penrose Model}
	
	In the Diósi-Penrose model, the rate reads:
	
	\begin{equation}
		\Gamma_{\text{DP}} = \frac{G M^2}{\hbar R},
	\end{equation}
	
	where \(R\) is the object size. For typical parameters (\(T_{\text{coh}} = 10\) s, \(R = 50\) nm), this yields \(M_{\max} \approx 5 \times 10^{10}\) u.
	
	\begin{center}
		\small
		\resizebox{\textwidth}{!}{%
			\begin{tabular}{p{0.45\textwidth}p{0.45\textwidth}}
				\textbf{Diósi-Penrose} & \textbf{FFGFT (T0)} \\
				\hline
				Heuristic & Structural from fractal vacuum geometry \\
				No fundamental scale & Scaling through \(\xi\) \\
				\(M_{\max} \propto \sqrt{R}\) & Logarithmic corrections through \(f(\Delta x)\) \\
				Parameter-dependent & Additional ξ contribution to decoherence \\
			\end{tabular}%
		}
	\end{center}
	
	The decisive difference: In the Diósi-Penrose model, gravitational decoherence is an ad-hoc postulate. In FFGFT, it \emph{emerges} from the fractal structure of the vacuum field — the phase gradients \(\partial_z \theta\) naturally couple to gravitational field differences.
	
	\subsection*{Experimental Perspective}
	
	Experiments such as MAQRO (space-based interferometry) and MAST-QG (laboratory) aim to separate gravitational decoherence from environmental decoherence. The FFGFT-specific contribution would be:
	
	\begin{itemize}
		\item \textbf{Logarithmic deviation:} The fractal correlation function produces a characteristic \(\ln(\Delta x)\) dependence that differs from the \(\sqrt{R}\) scaling of the DP model
		\item \textbf{ξ-signature:} A systematic deviation proportional to \(\xi \approx 1.33 \times 10^{-4}\) relative to standard decoherence models
	\end{itemize}
	
	The experimental challenge lies in isolating the small ξ-contribution from other decoherence sources.
	
	\subsection*{Conclusion}
	
	FFGFT provides a \emph{structural mechanism} for gravitational decoherence of macroscopic superpositions: the fractal nonlinearity of the vacuum field generates an additional decoherence source that scales with \(\xi\) and exhibits logarithmic corrections. A specific numerical prediction for \(M_{\max}\) requires the complete determination of the fractal correlation function and is the subject of further research. Experiments such as MAQRO and MAST-QG can in principle test the characteristic scaling signature of FFGFT.


